\documentclass[11pt]{article}

\usepackage[T1]{fontenc}
\usepackage{lmodern}
\usepackage{amsmath,amssymb}
\usepackage{hyperref}
\usepackage[margin=1in]{geometry}

\newcommand{\Sym}{\mathrm{Sym}}
\newcommand{\mc}[1]{\mathcal{#1}}
\newcommand{\bR}{\mathbb{R}}
\newcommand{\Aut}{\mathrm{Aut}}
\newcommand{\Inj}{\mathrm{Inj}}

\title{Finite-\texorpdfstring{$N$}{N} Asymptotic-State Normalization for General NS--NS Profiles}
\author{}
\date{}

\begin{document}
\maketitle

\section{Aim}
This note extends the companion single-cycle derivation \texttt{ASYMPTOTIC\_STATE\_MAP.tex} from one active long string on each tube to arbitrary NS--NS multi-string profiles on the three-punctured sphere
\[
C_{0,3}=S^2\setminus\{0,1,\infty\}.
\]
The goal is still purely kinematic: determine the exact finite-$N$ factor that converts a canonically labeled marked worldsheet amplitude into the physical amplitude of the symmetric-product orbifold
\begin{equation}\label{eq:aim-orbifold-general}
\Sym^N(\bR^8).
\end{equation}
I do \emph{not} derive a new BRST-complete local vertex operator for general profiles here. The only dynamical input is that the strict Poisson/BF worldsheet computation produces an $N$-independent marked fixed-cover weight, together with the already-derived local cylinder factor $\sqrt{w}$ for one constituent string of winding $w$.

The basic convention is that a profile
\begin{equation}\label{eq:active-profile}
\lambda=(w_1,\dots,w_{\ell(\lambda)}),
\qquad
w_a\in \mathbb{Z}_{>0},
\qquad
|\lambda|:=\sum_{a=1}^{\ell(\lambda)} w_a\le N,
\end{equation}
records only the \emph{active} constituent strings in the external state. The omitted $N-|\lambda|$ copies of the orbifold stay in the identity sector and play no dynamical role in the connected worldsheet diagram. Active winding-one strings may still appear explicitly inside $\lambda$ when they carry nontrivial internal state data.

\section{Decorated Profiles and Marked Tube States}
Fix one tube $z_i$ and one active profile $\lambda=(w_1,\dots,w_{\ell(\lambda)})$. To each constituent string attach a transverse NS--NS seed label $\chi_a$, and write
\begin{equation}\label{eq:decorated-profile-data}
\vec\chi:=(\chi_1,\dots,\chi_{\ell(\lambda)}).
\end{equation}
Two constituent strings are treated as genuinely identical only when both their winding numbers and their seed labels agree. The corresponding permutation group is
\begin{equation}\label{eq:identical-constituent-group}
G_{\lambda,\vec\chi}
:=
\left\{
\pi\in S_{\ell(\lambda)}
\ \middle|\
(w_{\pi(a)},\chi_{\pi(a)})=(w_a,\chi_a)\ \text{for all }a
\right\}.
\end{equation}
If all repeated windings carry the same seed data, then $|G_{\lambda,\vec\chi}|=\prod_n m_n(\lambda)!$, where $m_n(\lambda)$ is the multiplicity of the part $n$ in $\lambda$. If repeated windings carry distinct seed labels, then the corresponding factor in $G_{\lambda,\vec\chi}$ is absent.

For each constituent string, keep the marked local sheet set
\[
[w_a]_{(a)}:=\{1,\dots,w_a\}_{(a)}.
\]
The full ordered marked sheet set is the disjoint union
\begin{equation}\label{eq:marked-sheet-set}
B_\lambda
:=
\bigsqcup_{a=1}^{\ell(\lambda)} [w_a]_{(a)},
\qquad
|B_\lambda|=|\lambda|.
\end{equation}
An ordered embedding of this marked active data into the $N$ physical copy labels is an injection
\begin{equation}\label{eq:general-ordered-injection}
\iota\in \Inj(B_\lambda,[N]).
\end{equation}
As bookkeeping, let
\begin{equation}\label{eq:ordered-bookkeeping-state}
\bigl|\Psi^{\rm ord}_{\iota;\lambda,\vec\chi}\bigr\rangle_i
\end{equation}
denote the ordered asymptotic tube state associated with the injection $\iota$. Labels not in $\iota(B_\lambda)$ remain in the identity sector.

The group $G_{\lambda,\vec\chi}$ acts on ordered injections by permuting genuinely identical constituent strings:
\begin{equation}\label{eq:group-action-on-injections}
(\pi\cdot \iota)(x):=\iota(\pi^{-1}x),
\qquad
\pi\in G_{\lambda,\vec\chi},
\quad
x\in B_\lambda.
\end{equation}
Because the constituent blocks in \eqref{eq:marked-sheet-set} are disjoint and $\iota$ is injective, this action is free. The physically marked state therefore depends only on the orbit
\begin{equation}\label{eq:orbit-of-injection}
[\iota]\in \Inj(B_\lambda,[N]) / G_{\lambda,\vec\chi},
\end{equation}
which I denote by
\begin{equation}\label{eq:marked-physical-state}
\bigl|\Psi^{\rm mark}_{[\iota];\lambda,\vec\chi}\bigr\rangle_i.
\end{equation}
The number of such orbits is
\begin{equation}\label{eq:number-of-marked-orbits}
\left|
\Inj(B_\lambda,[N]) / G_{\lambda,\vec\chi}
\right|
=
\frac{(N)_{|\lambda|}}{|G_{\lambda,\vec\chi}|},
\qquad
(N)_m:=\frac{N!}{(N-m)!}.
\end{equation}

Different orbit labels differ on at least one active strand assignment, so the cylinder pairing is
\begin{equation}\label{eq:general-marked-orthogonality}
\left\langle
\Psi^{\rm mark,\,out}_{[\iota];\lambda,\vec\chi}
\middle|
\Psi^{\rm mark,\,in}_{[\iota'];\lambda',\vec\chi'}
\right\rangle_{\rm cyl}
=
\delta_{\lambda,\lambda'}\,
\delta_{\vec\chi,\vec\chi'}\,
\delta_{[\iota],[\iota']}.
\end{equation}
The normalized finite-$N$ external state on the $i$-th tube is therefore
\begin{equation}\label{eq:general-external-state}
\bigl|\Psi^{\rm ext}_{\lambda,\vec\chi}(N)\bigr\rangle_i
:=
\sqrt{
\frac{|G_{\lambda,\vec\chi}|}{(N)_{|\lambda|}}
}
\sum_{[\iota]\in \Inj(B_\lambda,[N]) / G_{\lambda,\vec\chi}}
\bigl|\Psi^{\rm mark}_{[\iota];\lambda,\vec\chi}\bigr\rangle_i,
\end{equation}
and \eqref{eq:general-marked-orthogonality} implies
\begin{equation}\label{eq:general-external-state-unit-norm}
\left\langle
\Psi^{\rm ext,\,out}_{\lambda,\vec\chi}(N)
\middle|
\Psi^{\rm ext,\,in}_{\lambda',\vec\chi'}(N)
\right\rangle_{\rm cyl}
=
\delta_{\lambda,\lambda'}\,
\delta_{\vec\chi,\vec\chi'}.
\end{equation}

\section{Connected Three-Point Data}
For the three tubes of $C_{0,3}$, let the active external data be
\[
(\lambda_i,\vec\chi_i),
\qquad
\lambda_i=(w_{i,1},\dots,w_{i,\ell_i}),
\qquad
\ell_i:=\ell(\lambda_i),
\qquad
i=1,2,3.
\]
The genus-zero connected degree determined by these three active profiles is fixed by Riemann--Hurwitz:
\begin{equation}\label{eq:general-degree-d0}
d_0(\lambda_1,\lambda_2,\lambda_3)
:=
1+\frac12\sum_{i=1}^3\bigl(|\lambda_i|-\ell_i\bigr).
\end{equation}
If \eqref{eq:general-degree-d0} is not a nonnegative integer, or if no transitive monodromy triple with the required cycle data exists, then the connected genus-zero amplitude is defined to vanish.

The normalized connected three-point amplitude is
\begin{equation}\label{eq:general-external-amplitude}
\mc A_{123}^{\rm ext,(0)}(N)
:=
\left\langle
\Psi^{\rm ext,\,out}_{\lambda_3,\vec\chi_3}(N)
\middle|
\mc P_{0,3}
\middle|
\Psi^{\rm ext,\,in}_{\lambda_1,\vec\chi_1}(N)\otimes
\Psi^{\rm ext,\,in}_{\lambda_2,\vec\chi_2}(N)
\right\rangle.
\end{equation}
Expanding the three external averages \eqref{eq:general-external-state} gives
\begin{equation}\label{eq:general-external-amplitude-expanded}
\mc A_{123}^{\rm ext,(0)}(N)
=
\sqrt{
\frac{
|G_{\lambda_1,\vec\chi_1}|\,
|G_{\lambda_2,\vec\chi_2}|\,
|G_{\lambda_3,\vec\chi_3}|
}{
(N)_{|\lambda_1|}\,
(N)_{|\lambda_2|}\,
(N)_{|\lambda_3|}
}
}
\sum_{[\iota_1],[\iota_2],[\iota_3]}
\mc A_{123}^{\rm mark,(0)}([\iota_1],[\iota_2],[\iota_3]).
\end{equation}

The marked amplitude vanishes unless the three external profiles fit together into one connected degree-$d_0$ branched cover of the target sphere. Equivalently, on the canonical active sheet set $[d_0]$ one needs a monodromy triple
\begin{equation}\label{eq:general-connected-monodromy}
(h_1,h_2,h_3)\in S_{d_0}^3
\end{equation}
such that
\begin{itemize}
\item $h_i$ has cycle type $\lambda_i\,1^{d_0-|\lambda_i|}$,
\item $h_1h_2h_3=1$,
\item the subgroup $\langle h_1,h_2,h_3\rangle$ acts transitively on $[d_0]$.
\end{itemize}
Here the notation $\lambda_i\,1^{d_0-|\lambda_i|}$ means that the nontrivial cycles specified by $\lambda_i$ are present, while the remaining active sheets are fixed at the $i$-th puncture. Thus the external active profile does \emph{not} need to sum to $d_0$; the missing sheets are exactly the active copies that appear trivially at that puncture.

The weighted Hurwitz count is
\begin{equation}\label{eq:general-hurwitz-weighted}
H_{d_0}^{(0)}(\lambda_1,\lambda_2,\lambda_3)
:=
\sum_{[f]}
\frac{1}{|\Aut(f)|},
\end{equation}
where the sum runs over connected genus-zero degree-$d_0$ cover classes with the monodromy data described above.

\section{Counting Finite-\texorpdfstring{$N$}{N} Embeddings}
Fix one connected cover class $[f]$ of degree $d_0=d_0(\lambda_1,\lambda_2,\lambda_3)$, represented by a monodromy triple \eqref{eq:general-connected-monodromy} on the canonical active set $[d_0]$. As in the single-cycle derivation, simultaneous conjugation by $S_{d_0}$ gives an orbit of size
\begin{equation}\label{eq:general-orbit-size}
\frac{d_0!}{|\Aut(f)|}.
\end{equation}
To realize this active-sheet cover inside the finite-$N$ physical labels, choose an ordered injection
\begin{equation}\label{eq:general-cover-injection}
\jmath\in \Inj([d_0],[N]).
\end{equation}
The resulting concrete triple on $[N]$ acts by $\jmath h_i\jmath^{-1}$ on $\jmath([d_0])$ and trivially on the complement.

The local marked puncture data and the quotient by $G_{\lambda_i,\vec\chi_i}$ are already encoded in the external states \eqref{eq:general-external-state}. They do \emph{not} modify the global active-sheet multiplicity \eqref{eq:general-orbit-size}. Therefore the number of concrete finite-$N$ connected realizations associated with the cover class $[f]$ is still
\begin{equation}\label{eq:general-concrete-triple-count}
\frac{(N)_{d_0}}{|\Aut(f)|}.
\end{equation}

Let
\begin{equation}\label{eq:general-marked-cover-weight}
\mc W_f^{\rm mark}(\lambda_1,\vec\chi_1;\lambda_2,\vec\chi_2;\lambda_3,\vec\chi_3)
\end{equation}
denote the marked fixed-cover weight of one canonically labeled active-sheet realization. By construction this quantity is independent of $N$. Combining \eqref{eq:general-external-amplitude-expanded} with \eqref{eq:general-concrete-triple-count} gives
\begin{equation}\label{eq:general-external-amplitude-factorized}
\mc A_{123}^{\rm ext,(0)}(N)
=
\mc G_N^{\rm ext}(\lambda_1,\vec\chi_1;\lambda_2,\vec\chi_2;\lambda_3,\vec\chi_3)
\sum_{[f]}
\mc W_f^{\rm mark}(\lambda_1,\vec\chi_1;\lambda_2,\vec\chi_2;\lambda_3,\vec\chi_3),
\end{equation}
with the exact external-state factor
\begin{equation}\label{eq:general-derived-external-factor}
\mc G_N^{\rm ext}(\lambda_1,\vec\chi_1;\lambda_2,\vec\chi_2;\lambda_3,\vec\chi_3)
:=
\frac{(N)_{d_0}\,
\sqrt{
|G_{\lambda_1,\vec\chi_1}|\,
|G_{\lambda_2,\vec\chi_2}|\,
|G_{\lambda_3,\vec\chi_3}|
}
}{
\sqrt{
(N)_{|\lambda_1|}\,
(N)_{|\lambda_2|}\,
(N)_{|\lambda_3|}
}
}.
\end{equation}

The already-derived single-string cylinder analysis suggests that the marked fixed-cover weight factorizes as a product of local constituent-string normalizations,
\begin{equation}\label{eq:general-local-factorization-input}
\mc W_f^{\rm mark}
=
\left(
\prod_{i=1}^3
\prod_{a=1}^{\ell_i}
\sqrt{w_{i,a}}
\right)
\widetilde{\mc W}_f,
\end{equation}
where $\widetilde{\mc W}_f$ is the remaining reduced fixed-cover weight. Equation \eqref{eq:general-local-factorization-input} is imported input here, not a new derivation. Substituting it into \eqref{eq:general-external-amplitude-factorized} gives
\begin{equation}\label{eq:general-external-amplitude-local-factor}
\mc A_{123}^{\rm ext,(0)}(N)
=
\mc G_N^{\rm ext}
\left(
\prod_{i=1}^3
\prod_{a=1}^{\ell_i}
\sqrt{w_{i,a}}
\right)
\sum_{[f]}
\widetilde{\mc W}_f.
\end{equation}

\section{Relation to Decorated Class Normalization}
The orbifold operator that matches the active profile $\lambda=(w_1,\dots,w_{\ell(\lambda)})$ with seed labels $\vec\chi$ is normalized over \emph{unmarked} constituent cycles. Starting from ordered injections of the marked sheet set $B_\lambda$, there are two quotients:
\begin{itemize}
\item each constituent $w_a$-cycle has a local cyclic redundancy of order $w_a$,
\item the subgroup $G_{\lambda,\vec\chi}$ permutes genuinely identical decorated constituent strings.
\end{itemize}
Therefore the size of the decorated finite-$N$ class is
\begin{equation}\label{eq:general-decorated-class-size}
|[\lambda,\vec\chi]|_N
=
\frac{(N)_{|\lambda|}}{
\left(\prod_{a=1}^{\ell(\lambda)} w_a\right)\,
|G_{\lambda,\vec\chi}|
}
=
\frac{N!}{
(N-|\lambda|)!\,
\left(\prod_{a=1}^{\ell(\lambda)} w_a\right)\,
|G_{\lambda,\vec\chi}|
}.
\end{equation}
If the seed labels are suppressed and identical constituent strings are identified only by their winding numbers, then \eqref{eq:general-decorated-class-size} reduces to
\begin{equation}\label{eq:general-undecorated-class-size}
|[\lambda]|_N
=
\frac{N!}{
(N-|\lambda|)!\,
\prod_n n^{m_n(\lambda)}\,m_n(\lambda)!
}.
\end{equation}

Combining \eqref{eq:general-derived-external-factor} with the local product in \eqref{eq:general-local-factorization-input} yields the exact identity
\begin{equation}\label{eq:general-class-normalization-identity}
\mc G_N^{\rm ext}(\lambda_1,\vec\chi_1;\lambda_2,\vec\chi_2;\lambda_3,\vec\chi_3)
\left(
\prod_{i=1}^3
\prod_{a=1}^{\ell_i}
\sqrt{w_{i,a}}
\right)
=
\frac{N!}{(N-d_0)!}
\prod_{i=1}^3
\frac{1}{\sqrt{|[\lambda_i,\vec\chi_i]|_N}}.
\end{equation}
Equation \eqref{eq:general-class-normalization-identity} is the general NS--NS version of the single-cycle identity in \texttt{ASYMPTOTIC\_STATE\_MAP.tex}. It is purely a normalization statement: no orbifold three-point coefficient has been fitted anywhere in the derivation.

\section{Checks}
The reduction to the previous note is immediate. If each profile is single-cycle,
\[
\lambda_i=(w_i),
\qquad
|G_{\lambda_i,\vec\chi_i}|=1,
\qquad
|\lambda_i|=w_i,
\qquad
\ell_i=1,
\]
then \eqref{eq:general-degree-d0} gives
\[
d_0=\frac{w_1+w_2+w_3-1}{2},
\]
and \eqref{eq:general-derived-external-factor} reproduces the exact factor in \texttt{ASYMPTOTIC\_STATE\_MAP.tex}.

For one tube with profile $\lambda=(2,2)$ and identical seed data $\vec\chi=(\chi,\chi)$, the identical-constituent group has size
\[
|G_{\lambda,\vec\chi}|=2!,
\]
so \eqref{eq:general-decorated-class-size} gives
\[
|[(2,2),(\chi,\chi)]|_N
=
\frac{N!}{2!\,2^2\,(N-4)!}.
\]
If instead the two constituents have distinct seed labels, $\vec\chi=(\chi_1,\chi_2)$ with $\chi_1\neq \chi_2$, then $|G_{\lambda,\vec\chi}|=1$ and
\[
|[(2,2),(\chi_1,\chi_2)]|_N
=
\frac{N!}{2^2\,(N-4)!}.
\]
The extra factor of $2!$ is therefore exactly the quotient by identical-particle symmetry and nothing else.

Active winding-one strings are included explicitly in the same way. For example, the profile $\lambda=(1,1,2)$ has $|\lambda|=4$ and $\ell(\lambda)=3$. The two winding-one constituents contribute to \eqref{eq:general-degree-d0}, to the orbit count \eqref{eq:number-of-marked-orbits}, and to the class size \eqref{eq:general-decorated-class-size}. They are not part of the omitted passive spectators counted by $N-|\lambda|$.

Finally, with fixed profiles and $N\to\infty$, equation \eqref{eq:general-derived-external-factor} expands as
\begin{equation}\label{eq:general-large-n-expansion}
\mc G_N^{\rm ext}
=
N^{\,1-\frac12(\ell_1+\ell_2+\ell_3)}
\sqrt{
|G_{\lambda_1,\vec\chi_1}|\,
|G_{\lambda_2,\vec\chi_2}|\,
|G_{\lambda_3,\vec\chi_3}|
}
\left[
1+\frac{1}{N}
\left(
\frac14\sum_{i=1}^3 |\lambda_i|\bigl(|\lambda_i|-1\bigr)
-\frac12 d_0(d_0-1)
\right)
+O(N^{-2})
\right].
\end{equation}
Since
\[
d_0-\frac12\sum_{i=1}^3 |\lambda_i|
=
1-\frac12\sum_{i=1}^3 \ell_i,
\]
the leading power of $N$ depends only on the number of constituent strings on the three tubes. For $\ell_1=\ell_2=\ell_3=1$, this reduces to the previously derived $N^{-1/2}$ scaling.

\end{document}
